\section{Signale und Systeme}
\label{sec:signals_and_systems}
In diesem Abschnitt werden einige für den weiteren Verlauf wichtige Begriffe definiert, insbesondere \textit{Signal} und \textit{System},
sowie die sogenannten Linearen Zeit-Invarianten Systeme.
\paragraph{Signale}
Ein Signal beschreibt, wie ein Parameter sich im Verhältnis zu einem anderen Parameter verändert.
Signale übertragen Informationen über den Zustand oder das Verhalten eines physischen Systems. Mathematisch sind Signale
Funktionen einer oder mehrerer unabhängiger Variablen. Die unabhängige Variable von Audiosignalen ist Zeit oder Frequenz \cite[Kap. 5]{smith1997}.
Die unabhängige Variable kann stetig
oder diskret sein, wobei stetige Signale oft als \textit{analoge} Signale bezeichnet werden \cite[Kap. 2]{oppenheim1999}. In Computergestützten Anwendungen werden 
ausschließlich diskrete Signale verarbeitet.
Der Definitionsbereich Diskreter Signale sind die natürlichen Zahlen. Diskrete Signale sind also mathematisch betrachtet Reihen. 
Signale $S:\mathbb{N}\mapsto\mathbb{Z}$ heißen \textit{digitale Signale} \cite[Kap. 2]{oppenheim1999}.
Stetige Signale werden mit der unabhängigen Variable in runden Klammern notiert:
$x(t)$ für ein stetiges Signal $x$ in Abhängigkeit von Zeit $t$.
Diskrete Signale werden mit der unabhängigen Variable in eckigen Klammern notiert: $x[n]$ für ein diskretes Signal $x$,
wobei $n$ den Index des betrachteten Funktionswerts bezeichnet \cite[Seite 9]{oppenheim1999}.
\paragraph{Der Einheitsimpuls}
Ein besonderes diskretes Signal ist die \textit{Einheits-Sample-Reihe}.
Sie ist definiert als die Reihe
\begin{equation}
    \delta[n]=
    \begin{cases}
        0,\,n\ne0\\
        1,\,n=0
    \end{cases}
\end{equation}
Jedes diskrete Signal lässt sich als eine Summe von skalierten und verzögerten Einheitsimpulsen darstellen \cite[Kap. 2, Seite 11]{oppenheim1999}.
Im weiteren Verlauf wird die Einheits-Sample-Reihe etwas vereinfacht als \textit{Einheitsimpuls} bezeichnet.
\paragraph{Systeme}
Ein System produziert als Antwort auf ein Eingangssignal ein Ausgangssignal \cite[Kap. 5]{smith1997}.
Formal definiert \cite[Seite 16]{oppenheim1999} ein diskretes System als eine Transformation oder einen Operator, welche(r) ein Eingangssignal $x[n]$
auf ein Ausgangssignal $y[n]$ abbildet. Die Funktion $T:x[n]\mapsto y[n]$ heißt \textit{Transferfunktion} des Systems.
\paragraph{Lineare Zeit-Invariante Systeme (LTI)}
Besonders signifikant sind in der Signalverarbeitung Lineare Zeit-Invariante Systeme. Damit ein System \textit{linear}
heißt, muss es die beiden Eigenschaften
\begin{equation}
    T(x_1[n]+x_2[n])=T(x_1[n])+T(x_2[n])=y_1[n]+y_2[n]
\end{equation}
und 
\begin{equation}
    T(a\cdot x[n])=a\cdot T(x[n])=a\cdot y[n]
\end{equation}
erfüllen.
Daraus resultiert die \textit{Superposition} von Eingangssignalen:
\begin{equation}
    x[n]=\sum_{k}a_kx_k[n]\Rightarrow y[n]=\sum_{k}a_ky_k[n]
\end{equation}
Wobei $x[n]$ das Eingangssignal, $y[n]$ das Ausgangssignal, und $a_k$ eine Reihe von Konstanten bezeichnen \cite{oppenheim1999}.

Ein System heißt \textit{Zeit-Invariant}, wenn für jedes $n_0$ und jedes Eingangssignal $x$ mit den Werten
$x'[n]=x[n-n_0]$ das entsprechende Ausgangssignal $y$ mit den Werten $y'[n]=y[n-n_0]$ erzeugt wird \cite{oppenheim1999}.

Ist ein System sowohl linear als auch Zeit-Invarant, dann heißt es \textit{lineares Zeit-Invariantes System (LTI-System)}.
Die Auswirkungen von LTI-Systeme auf ein beliebiges Eingangssignal können vollständig durch ihre Impulsantwort (engl.: \textit{Impulse Response}, kurz \textit{IR},
siehe \ref{sec:convolution}) ausgedrückt werden \cite{oppenheim1999}.
Die Impulsantwort eines Systems ist selbst ein diskretes Signal und wird mit $h[n]$ notiert \cite[Kapitel 2, Seite 23]{oppenheim1999}.
